%% Präambel

% -- Dokumentenklasse
\documentclass[11pt,a4paper,oneside,ngerman]{article}

% -- Meta-Daten
\newcommand{\mytitle}{Der Titel meiner Arbeit}
\newcommand{\myname}{Jan Doant}

% -- Seitenränder
\usepackage[left=30mm,right=20mm,top=25mm,bottom=20mm]{geometry}

% -- Zeilenabstand
\usepackage[onehalfspacing]{setspace}

% -- keine Einrückung der ersten Zeile eines neuen Absatzes 
\setlength{\parindent}{0em}

%% -- Gestaltung Kopf- und Fußzeile
% -- Kopfzeile
\usepackage{fancyhdr}
\addtolength{\headheight}{3pt}
\lhead{\mytitle}
% -- Fußzeile
\lfoot{\myname}
\cfoot{}
\rfoot{Seite~\thepage}  

%% -- deutsches Sprachpaket
\usepackage[ngerman]{babel}
% -- Umlaute
\usepackage[utf8]{inputenc}
\usepackage[T1]{fontenc}
\usepackage{ae,aecompl}
% -- deutsche Anführungszeichen
\usepackage[babel,german=quotes]{csquotes}

% -- Mathematikpakete
\usepackage{amsfonts}
\usepackage{amsmath}

% -- Grafikpaket
\usepackage{graphicx}

% -- Verweise werden im Dokument klickbar
\usepackage{hyperref}
\usepackage[dvipsnames,svgnames,x11names,hyperref]{xcolor}
\hypersetup{
	pdfauthor={\myname},
	pdftitle={\mytitle},
	colorlinks = true, 
	linkcolor=[RGB]{63 81 181},
    citecolor= [RGB]{76 175 80},
    urlcolor=[RGB]{0 76 153},
}

% -- Verweise automatisieren
%\usepackage{cleveref}
\usepackage{varioref}
\usepackage{prettyref}
% -- prettyref-Befehl konfigurieren
\newrefformat{fig}{Abbildung \vref{#1}}
\newrefformat{tab}{Tabelle \vref{#1}}

% -- Anzeigen der verwendeten Labels (nur im Dev-Modus verwenden)
\usepackage[inner, 			% Information wird am inneren Rand angezeigt 
			%final			% blendet Informationen aus, wenn aktiv 
			]{showlabels}

% -- Bibliografiestil
\usepackage[
	sorting=none, 		%Literaturverzeichnis wird in der Reihenfolge, in der die Referenzen genannt werden
						% sortiert isbn=false,
	url=false, 			%keine Angabe der URL
	style=numeric,
	isbn=false,
	date=year,
	maxbibnames=6,      %maximal 6 Namen in der Literaturliste 
	minbibnames=3,      %mindestens 3 Namen in der Literaturliste 
	backend=biber]{biblatex}
% -- Pfad zur .bib-Datei	
\addbibresource{bib/bib.bib} 

%% Dokument
\begin{document}

% -- Titelseite
\pagestyle{empty}
\begin{titlepage}
	
\begin{figure}[htbp]
	\begin{minipage}{0.19\textwidth}
		% \textwidth bezieht sich nun auf die Minipage
		\includegraphics[width=\textwidth]{img/logo_tuc.pdf}
	\end{minipage}
	% Auffüllen des Zwischenraums
	\hfill
	% minipage mit Text
	\begin{minipage}{0.6\textwidth} 
		\begin{center}	
			Technische Universität Chemnitz\\
			Fakultät Maschinenbau\\
			Professur Fertigungsmesstechnik
		\end{center}
	\end{minipage}
	\hfill
	\begin{minipage}{0.19\textwidth} 
		% \textwidth bezieht sich nun auf die Minipage
		\includegraphics[width=\textwidth]{img/logo_fmt.pdf}
	\end{minipage}
\end{figure}
	
	
	\vspace*{1cm}
		
		\begin{center}	
		
		\vspace{1.5cm}
		
		{\huge\textbf{Bachelorarbeit}}
		
		\vspace{1.5cm}
		
		{\large Arbeitstitel:}\\
		{\large\textbf{\mytitle}}
		
		\vspace{1.5cm}
		
		Betreuer: M.Sc. Robert Hofmann\\
		Prüfer: Univ.-Prof. Dr.-Ing. Sophie Gröger		
		
		\vfill
		
			
		
	\end{center}
	
	\begin{flushleft}
		Name: Jan Doant\\		
		E-Mail: jan.doant@s2017.tu-chemnitz.de\\
		Studiengang: Bachelor Maschinenbau (7. Semester)\\
		Matrikelnummer: 461311\linebreak
		
		
		Eingereicht am: 06.06.2018
		
	\end{flushleft}
	
	
\end{titlepage}
\newpage

% -- Selbstständigkeitserklärung
\pagestyle{empty}
Selbststaendigkeitserklärung
\newpage

% -- Inhaltsverzeichnis
\pagestyle{empty}
%%Anzahl der maximal aufgeführten Gliederungspunkte
\setcounter{tocdepth}{3}
\tableofcontents
\newpage

% --Sandbox
\pagestyle{fancy}
\section{Sandbox}

\subsection{Sandbox--Zitierung}

Das ist zu eine zu zitierende Textstelle \cite{Klocke2008}.

\subsection{Sandbox--Längere Abschnitte}
Jemand musste Josef K. verleumdet haben, denn ohne dass er etwas Böses getan hätte, wurde er eines
Morgens verhaftet. Wie ein Hund! sagte er, es war, als sollte die Scham ihn überleben.

Als Gregor Samsa eines Morgens aus unruhigen Träumen erwachte, fand er sich in seinem Bett zu einem
ungeheueren Ungeziefer verwandelt. Und es war ihnen wie eine Bestätigung ihrer neuen Träume und
guten Absichten, als am Ziele ihrer Fahrt die Tochter als erste sich erhob und ihren jungen Körper
dehnte.
»Es ist ein eigentümlicher Apparat«, sagte der Offizier zu dem Forschungsreisenden und überblickte
mit einem gewissermaßen bewundernden Blick den ihm doch wohlbekannten Apparat. Sie hätten noch ins
Boot springen können, aber der Reisende hob ein schweres, geknotetes Tau vom Boden, drohte ihnen
damit und hielt sie dadurch von dem Sprunge ab.

In den letzten Jahrzehnten ist das Interesse an Hungerkünstlern sehr zurückgegangen. Aber sie
überwanden sich, umdrängten den Käfig und wollten sich gar nicht fortrühren. Jemand musste Josef K.
verleumdet haben, denn ohne dass er etwas Böses getan hätte, wurde er eines Morgens verhaftet. Wie
ein Hund! sagte er, es war, als sollte die Scham ihn überleben. Als Gregor Samsa eines Morgens aus
unruhigen Träumen erwachte, fand er sich



\subsection{Sandbox--Grafiken einbinden}

\begin{figure}[htbp] 
  \centering
     \includegraphics[width=0.7\textwidth]{img/logo_tuc.pdf}
  \caption{Logo der Technischen Universität Chemnitz}
  \label{fig:Bild1}
\end{figure}

\newpage

\newpage


\subsection{Sandbox--Referenzierung}

Ich beziehe mich auf \prettyref{fig:Bild1}.
\newpage

% -- Einführung
\pagestyle{fancy}
\section{Einführung}
\label{sec:einfuehrung}









\newpage

% -- Abbildungsverzeichnis
\pagestyle{fancy}
\listoffigures
\newpage

% -- Literaturverzeichnis
\pagestyle{fancy}
\printbibliography

\end{document}
